\documentclass[10pt]{article}
\usepackage[a4paper,margin=0.6in]{geometry}
\usepackage{amsmath,amssymb,amsfonts,mathtools}
\usepackage[dvipsnames]{xcolor}
\usepackage{array,booktabs,float}
\usepackage[colorlinks=true,linkcolor=blue]{hyperref}
\allowdisplaybreaks
\numberwithin{equation}{section}
\setlength{\jot}{4pt}
\setlength{\emergencystretch}{3em}
\newcommand{\Dh}{\hat\Delta}
\newcommand{\OO}{\mathcal O}
\newcommand{\GG}{\mathcal G}
\newcommand{\CC}{\mathcal C}
\newcommand{\new}[1]{{\color{blue}#1}}
\newcommand{\pre}{\frac{g^{2}}{4\pi^{2}}\frac{1}{(2\pi)^{3}}\frac{g^{2}}{2\pi^{2}}\frac{N_{c}}{k^{+}}}
\newcommand{\XX}{\frac{(x-y)\cdot(x'-y')}{(x-y)^{2}(x'-y')^{2}}}
\newcommand{\XD}{\frac{(x-y)\cdot(y'-y)\,(x'-y')\cdot(y'-y)}{(x-y)^{2}(x'-y')^{2}(y'-y)^{2}}}
\newcommand{\ph}{e^{-ik\cdot(y'-y)\frac{1}{\xi}}}
\newcommand{\Ucol}{\big[U^{e'c}(y')-U^{e'c}(x')\big]\big[U^{ec}(y)-U^{ce}(x)\big]\rho^{e'}(x')\rho^{e}(x)}

\title{Group II, Row IV: the $d_\perp$ term and the plus prescription}
\date{}

\begin{document}
\maketitle
\noindent Your expressions are in black, new ones in \new{blue}. The term is the third line of Eq.~(fin) of \texttt{RowIV\_FF\_userform.pdf}.

%==========================================================================================
\section{The term}\label{sec:term}
\begin{align}
\frac{d^3N}{d^3k}\bigg|_{d_\perp}={}&\pre\int_{\frac{k^{+}}{\vee}}^{\frac{k^{+}}{k^{+}+\Lambda}}d\xi\,\frac{1-\xi}{\xi}\int_{x,y,x',y'}\ph\notag\\
&\times\Big[\XX-d_\perp\XD\Big]\notag\\
&\times\Ucol .
\label{eq:term}
\end{align}
\begin{itemize}
\item No pole multiplies this term, so $d_\perp=2-\epsilon$ can be set to $2$ here (the difference is $O(\epsilon)$). With $X\equiv\frac{x-y}{(x-y)^2}$, $X'\equiv\frac{x'-y'}{(x'-y')^2}$ and $\Dh\equiv\frac{y'-y}{|y'-y|}$, the bracket is
\begin{equation}
\new{X\cdot X'-2(X\cdot\Dh)(X'\cdot\Dh)=X^iX'^j\big(\delta^{ij}-2\Dh^i\Dh^j\big)}.
\end{equation}
The tensor $\delta^{ij}-2\Dh^i\Dh^j$ is traceless. This is what decides the answer below.
\item Write
\begin{equation}
\new{\begin{aligned}F(\xi)\equiv{}&\int_{x,y,x',y'}\ph\Big[\XX-2\XD\Big]\\&\times\Ucol,\end{aligned}}
\end{equation}
so that Eq.~\eqref{eq:term} is $\pre\int_{k^+/\vee}^{k^+/(k^++\Lambda)}d\xi\,\frac{1-\xi}{\xi}F(\xi)$.
\item With $\xi=\frac{k^+}{p^++k^+}$ (so $d\xi\,\frac{1-\xi}{\xi}\to dp^+\frac{p^+}{(p^++k^+)^2}$):
\begin{equation}
\new{\frac{d^3N}{d^3k}\bigg|_{d_\perp}=\pre\int_{\Lambda}^{\vee-k^+}dp^+\,\frac{p^+}{(p^++k^+)^2}\,F\Big(\frac{k^+}{p^++k^+}\Big)}.
\label{eq:pplus}
\end{equation}
\end{itemize}

%==========================================================================================
\section{Answer}\label{sec:answer}
\textbf{Yes, you may write a plus prescription, but it subtracts zero at both ends.} The term has no $\log\frac{\vee}{\Lambda}$, no $\log\frac{\vee}{k^+}$ and no $\log\frac{k^+}{\Lambda}$. It is finite and goes entirely into the finite part.
\begin{itemize}
\item \textbf{$\xi\to1$ ($p^+\to0$, the $\Lambda$ end).} The integrand is $\propto(1-\xi)$: in Eq.~\eqref{eq:pplus} it is $\frac{p^+}{k^{+2}}F(1)\to0$.
\begin{itemize}
\item There is no $\frac{1}{1-\xi}$ (no $\frac{1}{p^+}$), so there is nothing for $\big[\frac{\xi}{1-\xi}\big]_+$ or for your $\big[\frac{k^+}{p^+}\big]_+$ to regulate.
\item In your definition, $\frac{p^+}{P^2}F=\frac{k^+}{p^+}f(p^+)$ with $f(p^+)=\frac{p^{+2}}{k^+P^2}F$ and $f(0)=0$. Hence $\big[\frac{k^+}{p^+}\big]_+f=\frac{k^+}{p^+}f$, and the term $f(0)\log\frac{k^+}{\Lambda}$ it would generate is zero.
\end{itemize}
\item \textbf{$\xi\to0$ ($p^+\to\vee$, the $\vee$ end).} In Eq.~\eqref{eq:pplus} the weight is $\frac{1}{p^+}$, so there would be a $\log\vee$ if $F(0)\neq0$. But (Sec.~\ref{sec:small})
\begin{equation}
\new{\begin{aligned}&F(\xi)=\frac{8\pi\,\xi^2}{\mathbf k^2}\,\CC+O(\xi^3),\\
&\CC\equiv\int_{x,x',y}\frac{(x-y)\cdot\hat{\mathbf k}\,(x'-y)\cdot\hat{\mathbf k}-\frac12(x-y)\cdot(x'-y)}{(x-y)^2(x'-y)^2}\big[U^{e'c}(y)-U^{e'c}(x')\big]\big[U^{ec}(y)-U^{ce}(x)\big]\rho^{e'}(x')\rho^e(x).\end{aligned}}
\label{eq:small}
\end{equation}
\begin{itemize}
\item The integrand $\frac{1-\xi}{\xi}F(\xi)$ is $O(\xi)$, i.e.\ $O(1/p^{+3})$ in Eq.~\eqref{eq:pplus}.
\item A plus prescription at $\xi=0$, $\big[\frac1\xi\big]_+$ acting on $f(\xi)=(1-\xi)F(\xi)$, subtracts $f(0)=0$. The $\delta(\xi)$ term it would generate, $f(0)\log\frac{\vee}{k^+}$, is zero.
\end{itemize}
\end{itemize}
\paragraph{Result.} Both limits can be moved to $0$ and $1$, with power corrections only:
\begin{equation}
\boxed{\new{\begin{aligned}\frac{d^3N}{d^3k}\bigg|_{d_\perp}={}&\pre\int_{0}^{1}d\xi\,\frac{1-\xi}{\xi}\int_{x,y,x',y'}\ph\\
&\times\Big[\XX-2\XD\Big]\\
&\times\Ucol\\
&-\pre\Big[\frac{\Lambda^2}{2(k^++\Lambda)^2}\,F(1)+\frac{4\pi}{\mathbf k^2}\frac{k^{+2}}{\vee^2}\,\CC\Big]\\
&+O\Big(\frac{\Lambda^3}{k^{+3}},\frac{k^{+3}}{\vee^3}\Big).\end{aligned}}}
\label{eq:result}
\end{equation}
\begin{itemize}
\item $F(1)=\int_{x,y,x',y'}e^{-ik\cdot(y'-y)}\big[\dots\big]\Ucol$ is $F$ at $\xi=1$.
\item With $\Lambda\to0$ and $\vee\to\infty$ only the integral from 0 to 1 remains. It depends on neither $\Lambda$ nor $\vee$, so it is part of the finite (scale-independent) cross section.
\item It enters neither JIMWLK, nor BFKL, nor DGLAP (it has no pole).
\end{itemize}

%==========================================================================================
\section{Why \texorpdfstring{$F(\xi)=O(\xi^2)$}{F(xi) = O(xi^2)}}\label{sec:small}
With $\Delta\equiv y'-y$,
\begin{equation}
\new{\begin{aligned}&F(\xi)=\int d^2\Delta\,e^{-ik\cdot\Delta/\xi}\big(\delta^{ij}-2\Dh^i\Dh^j\big)\,\GG^{ij}(\Delta),\\
&\GG^{ij}(\Delta)\equiv\int_{x,x',y}\frac{(x-y)^i}{(x-y)^2}\frac{(x'-y-\Delta)^j}{(x'-y-\Delta)^2}\,\OO\big|_{y'=y+\Delta},\end{aligned}}
\end{equation}
with $\OO$ the colour factor of Eq.~\eqref{eq:term}. The parent momentum is $\mathbf q=\mathbf k/\xi$, so $\xi\to0$ is the Fourier transform at large $|\mathbf q|$. It is governed by $\Delta\to0$, where $\GG^{ij}(\Delta)=\GG^{ij}(0)+O(|\Delta|)$.
\begin{enumerate}
\item \textbf{Angular integral.} With $\Delta=r(\cos\phi,\sin\phi)$,
\begin{equation*}
\delta^{ij}-2\Dh^i\Dh^j=-\begin{pmatrix}\cos2\phi&\sin2\phi\\\sin2\phi&-\cos2\phi\end{pmatrix},\qquad\int_0^{2\pi}d\phi\,e^{-iqr\cos(\phi-\phi_q)}e^{2i\phi}=-2\pi J_2(qr)e^{2i\phi_q},
\end{equation*}
so
\begin{equation}
\new{\int_0^{2\pi}d\phi\,e^{-i\mathbf q\cdot\Delta}\big(\delta^{ij}-2\Dh^i\Dh^j\big)=4\pi\,J_2(qr)\Big(\hat q^i\hat q^j-\frac12\delta^{ij}\Big)}.
\end{equation}
\item \textbf{Radial integral.} $\int_0^\infty dr\,r\,J_2(qr)\,\GG^{ij}(0)\,\chi(r)=\frac{2}{q^2}\GG^{ij}(0)+O(q^{-3})$ for any smooth $\chi$ with $\chi(0)=1$, because $\int_0^\infty dt\,t\,J_2(t)=2$. The $O(|\Delta|)$ terms of $\GG$ give $O(q^{-3})$.
\item Hence $F=\frac{8\pi}{q^2}\big(\hat q^i\hat q^j-\frac12\delta^{ij}\big)\GG^{ij}(0)+O(q^{-3})$. With $q=|\mathbf k|/\xi$ and $\hat q=\hat{\mathbf k}$, this is Eq.~\eqref{eq:small}.
\end{enumerate}
\begin{itemize}
\item Only the traceless part of $\GG^{ij}(0)$ enters. The trace $\GG^{ii}(0)$ (which is sensitive to large $y$) drops out, so $\CC$ is finite.
\item The $1/q^2$ comes from the discontinuity of $\Dh^i\Dh^j$ at $\Delta=0$. It is a power, not a logarithm.
\end{itemize}

\paragraph{Contrast: where the $\log\frac{\vee}{k^+}$ of Row IV is.} It is in the first line of Eq.~(fin), not in this term.
\begin{itemize}
\item There the $\frac{1-\xi}{\xi}$ of $P_{gg}$ comes with $\frac{d\xi}{\xi^2}$, and the weight is $-\log\frac{(y'-y)^2\mu^2}{4}$.
\item Since $\int d^2\Delta\,e^{-i\mathbf q\cdot\Delta}(-\log\Delta^2)\,\GG^{ii}(\Delta)=\frac{4\pi}{q^2}\GG^{ii}(0)+\dots$, the integrand is $\frac{1-\xi}{\xi^3}\cdot\frac{4\pi\xi^2}{\mathbf k^2}\GG^{ii}(0)\to\frac{4\pi}{\mathbf k^2}\frac{\GG^{ii}(0)}{\xi}$.
\item This gives $\log\frac{\vee}{k^+}\,\frac{4\pi}{\mathbf k^2}\GG^{ii}(0)$. Here $y=y'$ ($=z$) and $\frac{4\pi}{\mathbf k^2}$ is the soft emission from $z$: this is the BFKL coefficient $P_{\rm II.IV}=2N_c\int K(w,w',z)K(x,x',z)\OO(z,z)$ of \texttt{LogVeeKplus\_BFKL.pdf}.
\item The difference: the BFKL piece has the trace $\GG^{ii}(0)$ with $\frac{d\xi}{\xi^3}$; the $d_\perp$ term has the traceless part with $\frac{d\xi}{\xi}$.
\item In $p^+$ language (\texttt{RowIV\_dperp.pdf}, Step 6): for $p^+\gg k^+$ only the $\delta_{mm'}\delta_{kk'}$ part of the polarization sum ($c_2$) survives. The $d_\perp$ term comes from $c_1+c_3=\frac{d_\perp p^+}{P^2}$, i.e.\ the weight $\xi(1-\xi)$, which after the $z$ integral falls like $1/p^{+3}$.
\end{itemize}

%==========================================================================================
\section{Remarks on the plus prescription}\label{sec:remarks}
\begin{itemize}
\item Any subtraction is an identity: $\int\frac{d\xi}{\xi}f=\int\frac{d\xi}{\xi}\big[f(\xi)-f(c)\big]+f(c)\log\frac{\xi_1}{\xi_0}$. The only natural choices, $c=0$ (at the $\vee$ end) or the $p^+\to0$ end, give $f(c)=0$, so nothing changes.
\item Do \emph{not} subtract the $\xi=1$ value at the $\xi\to0$ end, i.e.\ $\int_{\xi_0}^{\xi_1}\frac{d\xi}{\xi}\big[f(\xi)-f(1)\big]+f(1)\log\frac{\xi_1}{\xi_0}$. It is exact, but it creates the pair $\pm F(1)\log\frac{\xi_1}{\xi_0}=\pm F(1)\big[\log\frac{\vee}{k^+}-\log(1+\frac{\Lambda}{k^+})\big]$. The $+$ part would look like an evolution logarithm, but $F(1)$ belongs to neither BFKL nor JIMWLK; it cancels against the $-$ part hidden in the subtracted integral.
\item Do not confuse this with $\frac{\xi}{(1-\xi)_+}$ in the first line of Eq.~(fin). That one is needed: it regulates the soft end of $P_{gg}$, and its $\delta(1-\xi)$ term carries $\log\frac{\vee}{\Lambda}$ (JIMWLK). The present term has no $\frac{1}{1-\xi}$.
\end{itemize}

%==========================================================================================
\section{Numerical check}\label{sec:num}
Script \texttt{plus\_check.py}, results in \texttt{plus\_check.json}.
\paragraph{Configuration.}
\begin{itemize}
\item $SU(2)$ in the adjoint ($SO(3)$ rotations); $U(y)$ smooth, $U\to1$ outside $|y|\approx3$.
\item Two Gaussian colour charges. The colour of the second is chosen so that $\int_x[1-U(x)]^{ce}\rho^e(x)=0$; then $\int_xX(x-y)[U(y)-U(x)]\rho(x)$ falls like $1/y^2$.
\item $3072^2$ grid on a $40\times40$ box; $\GG^{ij}(\Delta)$ by FFT. A smooth window $e^{-\Delta^2/(2\cdot7^2)}$ only changes $F$ at $q\lesssim1/7$.
\item Normalization: $\GG^{ii}(0)=1$. Here $\GG^{ij}(0)=@@G0@@$.
\end{itemize}
\paragraph{(a) Calibration.} For $\GG^{ij}=M^{ij}e^{-\Delta^2}$ the grid transform agrees with the exact $4\pi M^{ij}(\hat q^i\hat q^j-\frac12\delta^{ij})\int_0^\infty dr\,r\,J_2(qr)e^{-r^2}$ to $\max q^2|\delta F|=@@CAL@@$ ($q\le40$).
\paragraph{(b) Large $q=|\mathbf k|/\xi$.} $q^2F$ against $8\pi(\hat q^i\hat q^j-\frac12\delta^{ij})\GG^{ij}(0)$ ($d_\perp$ term), and $q^2F_E$ ($F$ with $-\log\Delta^2$ and the trace, the BFKL piece) against $4\pi\GG^{ii}(0)=4\pi$:
\begin{center}\small
\begin{tabular}{@{}rrrrr|rr@{}}\toprule
direction $\phi_q$ & $q$ & $q^2F$ & prediction & ratio & $q^2F_E$ & $/4\pi$\\\midrule
@@ASYM@@
\bottomrule
\end{tabular}
\end{center}
The imaginary parts are below $10^{-3}$ of the real parts.

\paragraph{(c) The lower limit $\xi_0=k^+/\vee$.} $|\mathbf k|=@@K@@$, $\hat{\mathbf k}=\hat e_1$. Here $T(\xi_0)=\int_{\xi_0}^1d\xi\,\frac{1-\xi}{\xi}F$ (this term) and $T_E(\xi_0)=\int_{\xi_0}^1d\xi\,\frac{1-\xi}{\xi^3}F_E$ (the BFKL piece). $\delta T$ is the change from the previous row; the predictions follow from Eq.~\eqref{eq:small} and from $F_E\simeq\frac{4\pi\xi^2}{\mathbf k^2}\GG^{ii}(0)$.
\begin{center}\small
\begin{tabular}{@{}rrrr|rrr@{}}\toprule
$\xi_0$ & $T(\xi_0)$ & $\delta T$ & predicted $\delta T\propto\xi_0^2$ & $T_E(\xi_0)$ & $\delta T_E$ & predicted $\delta T_E\propto\log$\\\midrule
@@XI@@
\bottomrule
\end{tabular}
\end{center}
\begin{itemize}
\item $T$ changes by $\propto\xi_0^2$: each halving of $\xi_0$ adds a quarter of the previous change. It converges to $T(0)=@@T0@@$.
\item $T_E$ grows by the same amount, $\approx\frac{4\pi}{\mathbf k^2}\log2$, at each halving: this is $\log\frac{\vee}{k^+}$.
\end{itemize}

\paragraph{(d) The upper limit $\xi_1=\frac{k^+}{k^++\Lambda}$.} Strip $\int_{\xi_1}^1d\xi\frac{1-\xi}{\xi}F$ against $\frac{(1-\xi_1)^2}{2}F(1)$, with $F(1)=@@F1@@$:
\begin{center}\small
\begin{tabular}{@{}rrrr@{}}\toprule
$\Lambda/k^+$ & strip & $\frac{(1-\xi_1)^2}{2}F(1)$ & ratio\\\midrule
@@UP@@
\bottomrule
\end{tabular}
\end{center}
The ratio tends to 1 as $\Lambda/k^+\to0$, and the strip falls like $(\Lambda/k^+)^2$: no $\log\Lambda$.
\end{document}
